\documentclass[12pt]{amsart}
% Default margins are too wide all the way around. I reset them here
\setlength{\hoffset}{-.75in}
\setlength{\textwidth}{6.5in}
\setlength{\voffset}{-.5in}
\setlength{\textheight}{9.0in}
\usepackage{amsmath}
\usepackage{amssymb}
\usepackage{amsthm}
\usepackage{pgfplots}
\usepackage{enumerate}
\usepackage{tikz}
\usetikzlibrary{shadows,arrows,arrows.meta,shapes,positioning,calc}% Required for the Circle and Triangle arrow tips
\usepackage{setspace}
\usepackage{siunitx}
\usepackage{enumitem}
\usepackage{tcolorbox}
\usepackage{array}
\usepackage{microtype}
\usepackage{hyperref}

\theoremstyle{conjecture}
\newtheorem{conjecture}{Conjecture}

\theoremstyle{definition}
\newtheorem{definition}{Definition}

\theoremstyle{question}
\newtheorem{question}{Question}

\theoremstyle{theorem}
\newtheorem{theorem}{Theorem}

\newenvironment{solution}
{\begin{proof}[Solution]}
{\end{proof}}

\newcommand{\sectionrule}{%
    \vspace{0.3em}
    \hrule
    \vspace{0.8em}
}

\title{Chaos and Fractals: Patterns from Simple Rules}
\author{Hartwell Markham, Rob Tompkins, Rekik Ziku}
\date{\today}


\begin{document}

    \maketitle
    This project explores how simple repeated rules can produce intricate patterns, connecting the study of chaos with the
    geometry of fractals.
    Chaos concerns systems whose behavior can become difficult to predict because small differences in their starting
    conditions grow over time, while fractals describe shapes and sets with detail at increasingly small scales.
    These ideas often meet, but fractal structure does not by itself imply chaotic behavior.
    Using Chapter~11 of Strogatz's \emph{Nonlinear Dynamics and Chaos} as a foundation, we will consider how mathematicians
    describe the size and complexity of such patterns through countability, uncountability, and fractal dimension.
    The Cantor set provides a central example: removing intervals repeatedly leaves an uncountable collection of points with
    zero total length.
    Box dimension offers another way to understand its size by tracking how many small boxes are needed to cover it as the
    scale becomes finer.

    Two examples will connect these abstract ideas to the emergence of recognizable patterns: the chaos game and the Devil's
    staircase.
    In the familiar chaos game, repeatedly moving partway toward a randomly chosen vertex of a triangle produces the
    Sierpi\'nski triangle, showing how random choices can reveal an organized geometric structure.
    Barnsley and Vince's \emph{The chaos game on a general iterated function system} explains why this approach can recover
    an attracting set under broad conditions, extending beyond the usual setting in which every rule shrinks distances.
    Per Bak's \emph{The Devil's staircase} examines a different source of intricate organization: systems with competing
    rhythms that synchronize over ranges of conditions, creating a hierarchy of plateaus separated by a Cantor-like set.
    Together, these examples frame the project around a shared question: how does repeated evolution produce structure
    across scales?
    By comparing a fractal generated by random iteration with fractal organization among the conditions for synchronization,
    we will use visual examples and accessible geometric ideas to connect unpredictability, order, and the different
    meanings of mathematical size.

    \begin{thebibliography}{9}

        \bibitem{strogatz}
        Steven H. Strogatz.
        \emph{Nonlinear Dynamics and Chaos: With Applications to Physics, Biology, Chemistry, and Engineering}.
        2nd ed., CRC Press, 2018.
        Chapter~11, ``Fractals.''

        \bibitem{barnsley-vince}
        Michael F. Barnsley and Andrew Vince.
        ``The chaos game on a general iterated function system.''
        \emph{Ergodic Theory and Dynamical Systems}, 31(4), 1073--1079, 2011.
        Published online in 2010.
        DOI: \texttt{10.1017/S0143385710000428}.

        \bibitem{bak}
        Per Bak.
        ``The Devil's staircase.''
        \emph{Physics Today}, 39(12), 38--45, 1986.

    \end{thebibliography}

\end{document}